% Discussion.

\section{Discussion}\label{sec:discussion}

The paper's central result is Theorem~\ref{thm:dichotomy}: on a compact
featured-graph class of bounded size, the uniform closure of the
message-passing GNN class equals exactly the continuous $1$-WL
invariant functionals. Neither direction of the equality is new. The
contribution is to apply the equality to the polarisation measurement
literature. Doing so separates the published polarisation measures
into two classes, MPNN-computable and provably not, and it lets us
predict the output structure of any GNN-based polarisation estimator
from its architecture alone.

This chapter interprets the theoretical and empirical results for
the applied audience the paper is written for, states the checks the
audit procedure of Section~\ref{sec:emp-reddit} amounts to, records the
paper's limitations, and closes with the sentence the paper's title
is a compression of.

\subsection{What the dichotomy says}

For an applied audience, the substantive content of
Theorem~\ref{thm:dichotomy} is that a message-passing GNN's output on any
featured graph is a function of a specific summary statistic of the
graph: its $1$-WL colour histogram. Two featured graphs with the
same histogram get the same output, regardless of everything else
about them. Two featured graphs with different histograms may get
different outputs, but the differences run through histogram
distinctions.

The $1$-WL histogram is a computable object: colour refinement
(Definition~\ref{def:wl}) computes it in polynomial time. On a specific
dataset one can ask, empirically, what MPNN-computable statistics
the histogram tracks. Homophily is one such statistic
(Proposition~\ref{prop:computable-measures}). Attribute assortativity and
modularity of a fixed partition are others. These are the measures
that an MPNN estimator on a featured graph can approximate.

Closure statistics such as the local clustering coefficient
(Definition~\ref{def:clustering}) and the triangle density (Definition~\ref{def:triangles})
are not $1$-WL invariant. Spectral statistics such as the generalised
Euclidean distance (Definition~\ref{def:ged}) are not $1$-WL invariant. The
random walk controversy score (Definition~\ref{def:rwc}) is plausibly outside
the class as well, though the paper does not construct a witness
(Remark~\ref{rem:rwc}). These are measures a message-passing GNN estimator
cannot compute at any depth, width, aggregation, or training
procedure. The
dichotomy says which side of the line each measure lies on.

The witness pair of Section~\ref{sec:warmup} shows that the impossibility is
not vacuous: even on six-vertex graphs, there exist $1$-WL equivalent
pairs on which the closure-based and spectral measures differ by
constants. The ensemble-level entanglement of
Proposition~\ref{prop:entanglement} shows that the impossibility survives on
graphs of arbitrary size and holds pointwise on the two ensembles,
so the estimator's output does not merely fail to approximate
closure on adversarial pairs; it fails on families of graphs of
arbitrary size.

\subsection{Reading GNN-based polarisation results}

A working reader of the polarisation measurement literature is now
in a position to reinterpret the empirical claims of any paper that
uses a message-passing GNN as its polarisation estimator. The
reinterpretation follows three steps.

The first step is to identify what polarisation measure the paper
claims to compute. If the measure is edge homophily, assortativity,
or modularity of a fixed partition, the GNN estimator's output can
in principle track the measure; the paper's claim is defensible
subject to the usual checks on training and evaluation.

If the measure is closure-based or spectral, the GNN estimator's
output cannot track the measure in the sense of arbitrary
approximation. The paper's claim to compute the measure is not
compatible with the dichotomy of Theorem~\ref{thm:dichotomy}. What the
estimator's output is actually a function of is some functional in
$\mathcal{F}_{\wl}$; the audit procedure of Section~\ref{sec:emp-reddit}
identifies that functional's shape.

The third step is to run the audit. Fit a graph-blind regression on
first-order neighbourhood features and report the cross-validated
$R^2$ of the fit. Apply a $1$-WL preserving rewiring and report
how much the estimator's output moves. Compare the estimator's
numerical output across polarising and non-polarising datasets. If
the three checks together give a coherent story of graph-blind
reproducibility, closure invariance, and cross-topic
non-separation, the paper's original claim to measure structural
polarisation is not identifiable from the estimator's output.

\subsection{Recommendations for practice}

For a working researcher developing a new GNN-based polarisation
estimator, the practical implications of the results in this paper
are as follows.

First, the architecture choice determines the class of
polarisation measures the estimator can compute. A message-passing
architecture places the estimator in $\mathcal{F}_{\wl}$. If the
target measurement is closure or spectral in nature, message-passing
is the wrong architecture; higher-order GNNs, subgraph-count
architectures, or spectral estimators are more suitable.

Second, unsupervised link-reconstruction training (as in EchoGAE and
InfoVGAE) trains the estimator to predict edges, not to measure
polarisation. Any polarisation signal in the output is a
consequence of the link-reconstruction objective's alignment with
the label distribution, not of the architecture's ability to
measure polarisation directly. Supervised training against
polarisation labels changes what the estimator computes but does
not enlarge the class of measurable polarisation functionals.

Third, the audit procedure of Section~\ref{sec:emp-reddit} is a
low-effort check. Fitting a graph-blind ridge regression takes
seconds; running a rewiring test takes minutes. Running the audit
before publishing an estimator's polarisation scores would flag
cases where the estimator's output is dominated by first-order
homophily, letting the authors state their claim accurately.

\subsection{Limitations}

The ensembles of Proposition~\ref{prop:entanglement} are structurally uniform:
every connected component is a triangle or a $6$-cycle, and each
component is monochromatic. The impossibility survives on these
structured ensembles at every $n$, but the ensembles do not resemble
real interaction data. The obstacle analysis of
Section~\ref{sec:obstacle} enumerates weakenings that reduce structural
uniformity in exchange for weaker impossibility: the two-block
planted-partition model gives depth-$1$ $1$-WL total variation
$0.28$ between matched parameter points at a clustering gap of
$0.07$; the $K$-block SBM with matched first two moments reduces the
total variation to $\approx 0.06$ but shrinks the clustering gap to
$\approx 0.015$. In the dense graphon limit, message-passing GNNs
are universal for continuous cut-metric functionals
\citep{KerivenPeyre2020}, and the ensemble-level impossibility does
not extend. A stronger ensemble-level impossibility on realistic
interaction-graph distributions would need either $K$-block SBM
asymptotics with $K \to \infty$ or new graph-limit machinery
interpolating between Benjamini-Schramm and graphon convergence.

The empirical audit uses EchoGAE and its four released Twitter
datasets. InfoVGAE and TIMME lie in the same MPNN-computable
class and admit the same audit; we defer their audits to follow-up
work. The Waller-Anderson release covers ten thousand subreddits
but does not include the underlying interaction graph, so a
Reddit-native audit requires reconstructing an interaction graph
from Pushshift or ConvoKit dumps.

The audit procedure identifies what functional the estimator
computes, in the sense of showing that a small graph-blind
regression reproduces the estimator's output. It does not by itself
identify why the estimator's output correlates (when it does) with
human polarisation labels; that would require a further step
connecting the estimator's implicit functional to the label
generation process, which we do not attempt here.

\subsection{Open problems}

Three problems the paper does not solve are worth stating.

The first is whether a stronger ensemble-level entanglement holds on
random-graph ensembles that resemble real interaction data.
Section~\ref{sec:obstacle} enumerates the boundary of what is currently
achievable with SBM and $d$-regular constructions; a proof or
disproof of a stronger statement on realistic ensembles is open.

The second is a full characterisation of the exact class of
functionals a specific GNN architecture computes on a specific graph
distribution. Proposition~\ref{prop:egonet-factorisation} says the class lies
in $F_M(\Delta(\mathcal{R}_T))$ for some continuous $F_M$; nailing
down $F_M$ for a specific architecture and training distribution is
an approximation question we do not resolve.

The third is whether the audit procedure extends to supervised
GNN estimators trained end-to-end against polarisation labels. The
audit as stated operates on unsupervised estimators (EchoGAE, part
of InfoVGAE) where the training pipeline is decoupled from the
label distribution. A supervised estimator's implicit functional
is trained toward the label distribution, and the audit's
identification of the implicit functional would need to account for
this coupling.

\subsection{Conclusion}

When a message-passing GNN polarisation estimator on an interaction
graph agrees with human polarisation labels, the agreement is
evidence that the labels are aligned with a functional in
$\mathcal{F}_{\wl}$ on that dataset. The agreement is not evidence
that the graph has echo-chamber or spectral polarisation structure.
